\section{The original WZ1 slope and norm estimates}

For comparison, we record the statements in WZ1 Sections 6--7.
Extremality in this section has the WZ1 meaning of
Definition~\ref{def:wz1-extremal}. These statements use direction
separation and $L^{3/2}$ norms; WZ2 replaces them by the CWA and
projected-area statements above.

\begin{lemma}[A dense tube segment with small projection; WZ1, Lemma 6.2]
\label{lem:wz1-dense-segment}
\uses{def:ad-set,def:tube}
Let $0\leq\alpha<1$, $0<\delta\leq\rho\leq1$, and $\eps>0$.
Let $U$ be a segment of length $\sqrt\rho$ in a $\delta$-tube $T$,
with its full transverse thickness, and suppose
$F\subset U$ satisfies $|F|\geq\delta^\eps|U|$. If $v\in S^2$
and $v\cdot F$ is a $(\rho,\alpha,\delta^{-\eps})_1$-ADset, then
\[
                 |v\cdot\dir(T)|
                     \lesssim\delta^{-2\eps/(1-\alpha)}\sqrt\rho.
\]
\end{lemma}
\begin{proof}
A longitudinally dense subset of $U$ has a projection whose length
is controlled below by the angle of its axis to $v^\perp$.
The AD covering bound at scale $\rho$ bounds that projection above.
Comparing the two estimates and raising to the power $1/(1-\alpha)$
gives the conclusion.
\end{proof}

\begin{lemma}[Large derivative on a populated interval; WZ1, Lemma 6.3]
\label{lem:wz1-large-derivative}
\uses{prop:wz1-c2-grains}
For every $\eps>0$, choose $\eta$ and $\delta$ sufficiently small.
Suppose a WZ1 smooth-grain configuration with loss $\eta$ has
$\delta$-separated directions. If an interval $J\subset[-1,1]$
satisfies $\sqrt\delta\leq|J|\leq\delta^\eps$,
\[
 \sum_{T\in\T}|Y(T)\cap(\R^2\times J)|\geq\delta^\eta|J|,
\]
and $E_Y\cap(\R^2\times J)$ is covered by at most
$\delta^{-\eta}|J|^{-2+\sigma_3}$ cubes of side $|J|$, then
$|f'(z)|\geq|J|$ throughout $J$.
\end{lemma}
\begin{proof}
\uses{lem:wz1-dense-segment}
If the derivative is small at one point, the second derivative bound
keeps the global grain direction nearly constant across $J$. The
dense-segment lemma then confines many directions near one plane,
inconsistent with the direction count and the mass in the interval.
This is the directional predecessor of WZ2 Lemma 6.6.
\end{proof}

\begin{lemma}[A slab with comparable derivative; WZ1, Lemma 6.4]
\label{lem:wz1-comparable-derivative}
\uses{prop:wz1-c2-grains}
For every $\eps>0$ and scale ceiling, there is a WZ1 smooth-grain
configuration below that ceiling for which all shadings lie in
$[-1,1]^2\times[0,\delta^\eps]$, and a number $m\geq\delta^\eps$
such that
\[
             m\leq |f'(z)|\leq2m\quad(0\leq z\leq\delta^\eps).
\]
\end{lemma}
\begin{proof}
\uses{lem:wz1-large-derivative}
Choose a height interval retaining sufficient shaded mass and the
extremal spatial covering bound. Obtain the derivative lower bound
there, then localize and translate so that its magnitude is comparable
throughout one retained slab. Start with a smaller source loss to
retain the requested grain and extremality bounds.
\end{proof}

\begin{proposition}[Nonsingular WZ1 grains; WZ1, Proposition 6.1]
\label{prop:wz1-nonsingular}
\uses{def:wz1-extremal,def:plane-map,def:ad-set}
At arbitrarily small scales there are WZ1 $\eps$-extremal
configurations with global grains whose slope obeys
$1\leq|f'|\leq2$ and $|f''|\leq1/100$, together with local grains
and a $1$-Lipschitz plane map. The two AD covering bounds are those
of Proposition~\ref{prop:wz1-grains}, with the same retained shading.
The plane map is evaluated on that shaded union, as in
Remark~\ref{rem:plane-domain}.
\end{proposition}
\begin{proof}
\uses{lem:wz1-comparable-derivative,lem:mild-rescaling}
Rescale the selected height slab and the transverse coordinates.
The first derivative becomes comparable to one, and choosing the
slab sufficiently short makes the normalized second derivative small.
Transport the plane map and the local covering estimates with the
same affine change of variables.
\end{proof}

\begin{corollary}[A large projection norm; WZ1, Corollary 6.5]
\label{cor:wz1-large-norm}
\uses{def:twisted-projection,def:wz1-extremal}
For every $\eps>0$ and scale ceiling, there are a WZ1
$\eps$-extremal tube configuration below that ceiling and a slope
with $1\leq|f'|\leq2$, $|f''|\leq1/100$, such that
\[
 \left\|\sum_{T\in\T}\one_{\pi_f(T)}\right\|_{L^{3/2}(\R^2)}
                         \geq\delta^{-1-\sigma_3/3+\eps}.
\]
\end{corollary}
\begin{proof}
\uses{prop:wz1-nonsingular}
The global grains make the projected shaded union small, while the
sum of the projected shaded masses remains large. H\"older's
inequality bounds the $L^{3/2}$ norm below by this mass divided by
the one-third power of the support area. Enlarging each projected
shading to $\pi_f(T)$ can only increase the norm.
\end{proof}

\begin{proposition}[A small projection norm; WZ1, Proposition 7.1, normalized slope]
\label{prop:wz1-small-norm}
\uses{def:twisted-projection}
For every $\eps>0$, and all sufficiently small $\delta$, a family
of $\delta$-tubes with $\delta$-separated directions in WZ1's
bounded chart satisfies
\[
 \left\|\sum_{T\in\T}\one_{\pi_f(T)}\right\|_{L^{3/2}(\R^2)}
                           \leq\delta^{-1-\eps}
\]
whenever $f(0)=0$, $1\leq|f'|\leq2$, and $|f''|\leq1/100$.
\end{proposition}
\begin{proof}
\uses{thm:wz1-cinematic,lem:wz1-cinematic-family}
Split the four line parameters into approximately $\delta^{-1}$
windows of one slope parameter $c$. In each window, subtracting
$ct$ leaves a cinematic family. Direction separation supplies its
one-dimensional counting bound. Apply the cinematic estimate and
sum over the windows.
\end{proof}

\begin{remark}[An absolute normalization of the slope]
\label{rem:wz1-slope-offset}
The displayed paper statement of Proposition 7.1 specifies derivative
bounds but no bound on $f(0)$. Its uniform conclusion needs an
absolute normalization, which is present in the later cinematic
family construction and in the formal twisted-projection input.
Indeed, take a single vertical tube and $f(z)=A+z$. The derivative
conditions hold for every $A$, but the projected tube has area
comparable to $A\delta$ for large positive $A$, so its indicator has
$L^{3/2}$ norm comparable to $(A\delta)^{2/3}$. At a fixed scale this
is unbounded as $A\to\infty$. A fixed bound on $f(0)$, or the
normalization $f(0)=0$ used above, removes this obstruction.
\end{remark}

\begin{theorem}[The WZ1 cinematic estimate; WZ1, Theorem 7.2]
\label{thm:wz1-cinematic}
\uses{def:cinematic-family}
Let $\cF\subset C^2(I)$ be a fixed cinematic family on a compact
interval. For every $\eps>0$, at sufficiently small $\delta$, a
$\delta$-separated finite $F\subset\cF$ satisfying
\[
              \#(F\cap B)\leq\delta^{-\eps}(r/\delta)
\]
for every ball $B$ of diameter $r\geq\delta$ obeys
\[
               \left\|\sum_{g\in F}\one_{N_\delta(\graph g)}
                    \right\|_{3/2}\leq\delta^{-\eps}.
\]
For a threshold uniform over families, an absolute $C^2$ bound must
be fixed first, as explained in Remark~\ref{rem:pyz-uniformity}.
The completed Lean interface with a fixed counting constant is
Theorem~\ref{thm:cinematic-l32}.
\end{theorem}
\begin{proof}
\uses{thm:pyz}
Use the one-dimensional counting case of the cinematic maximal
estimate, localize the coordinate domain, and absorb fixed constants
with the fine-scale threshold.
\end{proof}

\begin{lemma}[The three-parameter cinematic family; WZ1, Lemma 7.3]
\label{lem:wz1-cinematic-family}
\uses{def:cinematic-family}
If $f(0)=0$, $1\leq|f'|\leq2$, and $|f''|\leq1/100$ on
$[0,1]$, the functions
\[
          g_{a,b,d}(t)=a+bf(t)+dtf(t),\qquad a,b,d\in[-1,1],
\]
form a cinematic family with constants independent of $f$.
They also have a uniform absolute $C^2$ bound.
\end{lemma}
\begin{proof}
The determinant of the two-jets of $1,f,tf$ is
$2(f')^2-ff''$, which is uniformly bounded away from zero because
$|f|\leq2$. Thus the jet at any point controls all three parameter
differences. This proves nondegeneracy and the equivalence of the
parameter and $C^2$ metrics, from which doubling follows.
\end{proof}
